\documentclass[10pt,a4paper]{amsart}
\usepackage{amsmath,amssymb,amsthm,graphicx}
% Retained preamble for the Stacks Project chapter {CH} (shared project preamble).
% The chapter file pinned-original-source.tex only carries \input{{preamble}} on its first line,
% so no preamble text lives inside the chapter itself.  The Stacks Project uses one single shared
% preamble for every chapter: the snapshot stages it three times (sources/stacks-chow/preamble.tex,
% sources/stacks-etale-cohomology/preamble.tex and focused-seeds/AGP/sources/stacks-snapshot-20261011/
% preamble.tex) and those copies are identical in size.  The block below is copied mechanically,
% line for line, from sources/stacks-chow/preamble.tex lines 14-18 and 154-217: the Xy-pic setup and
% the native theorem declarations together with the mathematical macros.  Nothing was added, removed
% or redefined; the wrapper class and the amsmath/amssymb/amsthm/graphicx base packages are supplied
% by the assembler.
% For commutative diagrams we use Xy-pic
% Chapter: stacks-spaces-more-morphisms.tex -- More on Morphisms of Spaces.
\usepackage[all]{xy}

% We use 2cell for 2-commutative diagrams.
\xyoption{2cell}

% Theorem environments.
%
\theoremstyle{plain}
\newtheorem{theorem}[subsection]{Theorem}
\newtheorem{proposition}[subsection]{Proposition}
\newtheorem{lemma}[subsection]{Lemma}

\theoremstyle{definition}
\newtheorem{definition}[subsection]{Definition}
\newtheorem{example}[subsection]{Example}
\newtheorem{exercise}[subsection]{Exercise}
\newtheorem{situation}[subsection]{Situation}

\theoremstyle{remark}
\newtheorem{remark}[subsection]{Remark}
\newtheorem{remarks}[subsection]{Remarks}

\numberwithin{equation}{subsection}

% Macros
%
\def\lim{\mathop{\mathrm{lim}}\nolimits}
\def\colim{\mathop{\mathrm{colim}}\nolimits}
\def\Spec{\mathop{\mathrm{Spec}}}
\def\Hom{\mathop{\mathrm{Hom}}\nolimits}
\def\Ext{\mathop{\mathrm{Ext}}\nolimits}
\def\SheafHom{\mathop{\mathcal{H}\!\mathit{om}}\nolimits}
\def\SheafExt{\mathop{\mathcal{E}\!\mathit{xt}}\nolimits}
\def\Sch{\mathit{Sch}}
\def\Mor{\mathop{\mathrm{Mor}}\nolimits}
\def\Ob{\mathop{\mathrm{Ob}}\nolimits}
\def\Sh{\mathop{\mathit{Sh}}\nolimits}
\def\NL{\mathop{N\!L}\nolimits}
\def\CH{\mathop{\mathrm{CH}}\nolimits}
\def\proetale{{pro\text{-}\acute{e}tale}}
\def\etale{{\acute{e}tale}}
\def\QCoh{\mathit{QCoh}}
\def\Ker{\mathop{\mathrm{Ker}}}
\def\Im{\mathop{\mathrm{Im}}}
\def\Coker{\mathop{\mathrm{Coker}}}
\def\Coim{\mathop{\mathrm{Coim}}}

% Boxtimes
%
\DeclareMathSymbol{\boxtimes}{\mathbin}{AMSa}{"02}

%
% Macros for moduli stacks/spaces
%
\def\QCohstack{\mathcal{QC}\!\mathit{oh}}
\def\Cohstack{\mathcal{C}\!\mathit{oh}}
\def\Spacesstack{\mathcal{S}\!\mathit{paces}}
\def\Quotfunctor{\mathrm{Quot}}
\def\Hilbfunctor{\mathrm{Hilb}}
\def\Curvesstack{\mathcal{C}\!\mathit{urves}}
\def\Polarizedstack{\mathcal{P}\!\mathit{olarized}}
\def\Complexesstack{\mathcal{C}\!\mathit{omplexes}}
% \Pic is the operator that assigns to X its picard group, usage \Pic(X)
% \Picardstack_{X/B} denotes the Picard stack of X over B
% \Picardfunctor_{X/B} denotes the Picard functor of X over B
\def\Pic{\mathop{\mathrm{Pic}}\nolimits}
\def\Picardstack{\mathcal{P}\!\mathit{ic}}
\def\Picardfunctor{\mathrm{Pic}}
\def\Deformationcategory{\mathcal{D}\!\mathit{ef}}
\usepackage{cleveref}
\crefname{originallemma}{original lemma}{original lemmas}
\Crefname{originallemma}{Original lemma}{Original lemmas}
\makeatletter
\newlabel{descent-lemma-compare-sites}{{070S}{source}{}{sourcepointer-1}{}}
\newlabel{descent-lemma-compare-sites@cref}{{[originallemma][0][]070S}{[1][0][]source}{}{}{}}
\makeatother
\crefname{originallemma}{original lemma}{original lemmas}
\Crefname{originallemma}{Original lemma}{Original lemmas}
\makeatletter
\newlabel{descent-remark-Zariski-site-space}{{03FG}{source}{}{sourcepointer-2}{}}
\newlabel{descent-remark-Zariski-site-space@cref}{{[originallemma][0][]03FG}{[1][0][]source}{}{}{}}
\makeatother
\crefname{originallemma}{original lemma}{original lemmas}
\Crefname{originallemma}{Original lemma}{Original lemmas}
\makeatletter
\newlabel{descent-remark-change-topologies-ringed}{{070R}{source}{}{sourcepointer-3}{}}
\newlabel{descent-remark-change-topologies-ringed@cref}{{[originallemma][0][]070R}{[1][0][]source}{}{}{}}
\makeatother
\crefname{originallemma}{original lemma}{original lemmas}
\Crefname{originallemma}{Original lemma}{Original lemmas}
\makeatletter
\newlabel{more-morphisms-lemma-sheaf-differentials-etale-localization}{{04BX}{source}{}{sourcepointer-4}{}}
\newlabel{more-morphisms-lemma-sheaf-differentials-etale-localization@cref}{{[originallemma][0][]04BX}{[1][0][]source}{}{}{}}
\makeatother
\crefname{originallemma}{original lemma}{original lemmas}
\Crefname{originallemma}{Original lemma}{Original lemmas}
\makeatletter
\newlabel{morphisms-definition-sheaf-differentials}{{01UQ}{source}{}{sourcepointer-5}{}}
\newlabel{morphisms-definition-sheaf-differentials@cref}{{[originallemma][0][]01UQ}{[1][0][]source}{}{}{}}
\makeatother
\crefname{originallemma}{original lemma}{original lemmas}
\Crefname{originallemma}{Original lemma}{Original lemmas}
\makeatletter
\newlabel{sites-modules-definition-module-differentials}{{04BN}{source}{}{sourcepointer-6}{}}
\newlabel{sites-modules-definition-module-differentials@cref}{{[originallemma][0][]04BN}{[1][0][]source}{}{}{}}
\makeatother
\title{Sheafifying the differentials of a morphism of schemes on the small etale site gives the module of differentials of the morphism of etale sites}
\author{The Stacks Project Authors}
\date{stacks-project:spaces-more-morphisms.tex:local-snapshot-20261011; GNU FDL 1.2 or later}
\begin{document}
\maketitle
\noindent\textit{Editorial note.} All delivered mathematical text is the author's own, copied from the pinned Stacks Project chapter file spaces-more-morphisms.tex; no statement or proof step was rewritten, reordered or removed.\par\medskip
\noindent\textit{Editorial note.} Preamble provenance, as required for this source: the first line of the pinned chapter file is \textbackslash{}input\{preamble\}, and the native preamble was not staged with the source, so no preamble text lives inside the chapter. The retained preamble used here is the project's single Stacks-wide preamble (the same 217-line / 8116-byte file that the already published chow and etale-cohomology cards use), copied mechanically line for line from sources/stacks-chow/preamble.tex lines 14-18 (Xy-pic setup) and 154-217 (theorem declarations and mathematical macros). Not one character was changed, no package was added, no macro was defined, and the excerpt serves only as preamble, never as a span coordinate.\par\medskip
\noindent\textit{Editorial note.} The same-chapter pointers are located at their \textbackslash{}label lines inside the pinned chapter file; the cross-chapter pointers are located in the project's own tag list staged in the snapshot (sources/stacks-chow/tags.txt), at the verbatim lines carrying their canonical tags.\par\medskip
\noindent\textit{Editorial note.} The proof is complete except for the last sentence: the author writes 'We omit the verification' of the two identities showing c'' and c mutually inverse on local sections (proof line 941). The construction of both maps and the adjunction argument are inside the counted span, so the omission is a routine tail, not a missing argument.\par\medskip
\noindent\textit{Allowed standard prerequisites (editorial).}\begin{itemize}
\item The sheaf of differentials is insensitive to etale localization. More on Morphisms, Lemma labelled more-morphisms-lemma-sheaf-differentials-etale-localization; invoked at lines 881-883. Gives h\textasciicircum{}*Omega\_\{X/Y\} = Omega\_\{U/Y\} for an etale h : U -> X, which identifies the values of (Omega\_\{X/Y\})\textasciicircum{}a.
\item Universal property of the module of differentials on a ringed site. Modules on Sites, Definition labelled sites-modules-definition-module-differentials; invoked at lines 891-898 and 912-913. Produces the universal O\_\{X\_etale\}-linear map c out of the derivation d\textasciicircum{}a.
\item Comparison of the small Zariski and small etale sites of a scheme. Descent, Remark labelled descent-remark-change-topologies-ringed and Lemma labelled descent-lemma-compare-sites; invoked at lines 920-927. Realises the restriction to the small Zariski site as the pushforward along id\_\{small,etale,Zar\}, so the map c' has an adjoint c''.
\item Sheaves on the small Zariski site of a scheme are sheaves on the scheme. Descent, Remark labelled descent-remark-Zariski-site-space; invoked at lines 905-907. Lets a derivation on the etale site be restricted to an ordinary Y-derivation on X.
\end{itemize}\medskip
\ifdefined\hypertarget\hypertarget{sourcepointer-1}{}\fi
\noindent\textit{Source reference.} Original lemma 070S: sources/stacks-chow/tags.txt, source lines 8966--8966. Navigation pointer for \textbackslash{}ref\{descent-lemma-compare-sites\} used inside the counted span (line 924); the argument that consumes it is retained in the span.\par\medskip
\ifdefined\hypertarget\hypertarget{sourcepointer-2}{}\fi
\noindent\textit{Source reference.} Original lemma 03FG: sources/stacks-chow/tags.txt, source lines 4491--4491. Navigation pointer for \textbackslash{}ref\{descent-remark-Zariski-site-space\} used inside the counted span (line 907); the argument that consumes it is retained in the span.\par\medskip
\ifdefined\hypertarget\hypertarget{sourcepointer-3}{}\fi
\noindent\textit{Source reference.} Original lemma 070R: sources/stacks-chow/tags.txt, source lines 8965--8965. Navigation pointer for \textbackslash{}ref\{descent-remark-change-topologies-ringed\} used inside the counted span (line 865); the argument that consumes it is retained in the span.\par\medskip
\ifdefined\hypertarget\hypertarget{sourcepointer-4}{}\fi
\noindent\textit{Source reference.} Original lemma 04BX: sources/stacks-chow/tags.txt, source lines 5669--5669. Navigation pointer for \textbackslash{}ref\{more-morphisms-lemma-sheaf-differentials-etale-localization\} used inside the counted span (line 883); the argument that consumes it is retained in the span.\par\medskip
\ifdefined\hypertarget\hypertarget{sourcepointer-5}{}\fi
\noindent\textit{Source reference.} Original lemma 01UQ: sources/stacks-chow/tags.txt, source lines 2419--2419. Navigation pointer for \textbackslash{}ref\{morphisms-definition-sheaf-differentials\} used inside the counted span (line 873); the argument that consumes it is retained in the span.\par\medskip
\ifdefined\hypertarget\hypertarget{sourcepointer-6}{}\fi
\noindent\textit{Source reference.} Original lemma 04BN: sources/stacks-chow/tags.txt, source lines 5658--5658. Navigation pointer for \textbackslash{}ref\{sites-modules-definition-module-differentials\} used inside the counted span (line 876); the argument that consumes it is retained in the span.\par\medskip
\begin{lemma}
\label{lemma-match-modules-differentials}
Let $f : X \to Y$ be a morphism of schemes. Let
$f_{small} : X_\etale \to Y_\etale$ be the associated
morphism of small \'etale sites, see
Descent, Remark \ref{descent-remark-change-topologies-ringed}.
Then there is a canonical isomorphism
$$
(\Omega_{X/Y})^a = \Omega_{X_\etale/Y_\etale}
$$
compatible with universal derivations. Here the first module
is the sheaf on $X_\etale$ associated
to the quasi-coherent $\mathcal{O}_X$-module $\Omega_{X/Y}$, see
Morphisms, Definition \ref{morphisms-definition-sheaf-differentials},
and the second module is the one from
Modules on Sites,
Definition \ref{sites-modules-definition-module-differentials}.
\end{lemma}

\begin{proof}
Let $h : U \to X$ be an \'etale morphism. In this case the natural map
$h^*\Omega_{X/Y} \to \Omega_{U/Y}$ is an isomorphism, see
More on Morphisms,
Lemma \ref{more-morphisms-lemma-sheaf-differentials-etale-localization}.
This means that there is a natural $\mathcal{O}_{Y_\etale}$-derivation
$$
\text{d}^a : \mathcal{O}_{X_\etale} \longrightarrow (\Omega_{X/Y})^a
$$
since we have just seen that the value of $(\Omega_{X/Y})^a$ on any object
$U$ of $X_\etale$ is canonically identified with
$\Gamma(U, \Omega_{U/Y})$. By the universal property of
$\text{d}_{X/Y} :
\mathcal{O}_{X_\etale}
\to
\Omega_{X_\etale/Y_\etale}$
there is a unique $\mathcal{O}_{X_\etale}$-linear map
$c : \Omega_{X_\etale/Y_\etale} \to (\Omega_{X/Y})^a$
such that
$\text{d}^a = c \circ \text{d}_{X/Y}$.

\medskip\noindent
Conversely, suppose that $\mathcal{F}$ is an
$\mathcal{O}_{X_\etale}$-module
and $D : \mathcal{O}_{X_\etale} \to \mathcal{F}$ is a
$\mathcal{O}_{Y_\etale}$-derivation. Then we can simply restrict
$D$ to the small Zariski site $X_{Zar}$ of $X$. Since sheaves on $X_{Zar}$
agree with sheaves on $X$, see
Descent, Remark \ref{descent-remark-Zariski-site-space},
we see that $D|_{X_{Zar}} : \mathcal{O}_X \to \mathcal{F}|_{X_{Zar}}$
is just a ``usual'' $Y$-derivation. Hence we obtain a map
$\psi : \Omega_{X/Y} \longrightarrow \mathcal{F}|_{X_{Zar}}$
such that $D|_{X_{Zar}} = \psi \circ \text{d}$. In particular, if we
apply this with $\mathcal{F} = \Omega_{X_\etale/Y_\etale}$
we obtain a map
$$
c' :
\Omega_{X/Y}
\longrightarrow
\Omega_{X_\etale/Y_\etale}|_{X_{Zar}}
$$
Consider the morphism of ringed sites
$\text{id}_{small, \etale, Zar} : X_\etale \to X_{Zar}$
discussed in
Descent, Remark \ref{descent-remark-change-topologies-ringed} and
Lemma \ref{descent-lemma-compare-sites}.
Since the restriction functor $\mathcal{F} \mapsto \mathcal{F}|_{X_{Zar}}$
is equal to $\text{id}_{small, \etale, Zar, *}$, since
$\text{id}_{small, \etale, Zar}^*$ is left adjoint to
$\text{id}_{small, \etale, Zar, *}$ and since
$(\Omega_{X/Y})^a = \text{id}_{small, \etale, Zar}^*\Omega_{X/Y}$
we see that $c'$ is adjoint to a map
$$
c'' :
(\Omega_{X/Y})^a
\longrightarrow
\Omega_{X_\etale/Y_\etale}.
$$
We claim that $c''$ and $c'$ are mutually inverse.
This claim finishes the proof of the lemma.
To see this it is enough to show that $c''(\text{d}(f)) = \text{d}_{X/Y}(f)$
and $c(\text{d}_{X/Y}(f)) = \text{d}(f)$ if $f$ is a local section of
$\mathcal{O}_X$ over an open of $X$. We omit the verification.
\end{proof}

\end{document}
